\chapter{Design \& Conceptual Model}

\todo{Relevance cycle of HEVNER: Stakeholders and requirements}

Should we add a section about HEVNER here? Or should we structure more of the chapter by the hefner thing like User
Target Group and Requirements Analysis?

\section{Starting Point and Motivation}

Repeat some things from the introduction. Further define the need for a combined and more comprehensive overview.

\section{User Target Group}

Target Users: Dashboard: policymakers, universities, companies Admin UI: Domain experts

\section{Requirement Analysis}

\subsection{Functional Requirements}

\subsubsection{Dashboard}
\label{FR1}

\subsubsection{Data Gathering Pipeline}
\label{FR2}

\subsubsection{Pipeline UI}
\label{FR3}

\subsection{Non-Functional Requirements}

\subsubsection{Maintainability and Code Quality}
\label{NFR1}

\subsubsection{Data Quality}
\label{NFR2}

\subsubsection{Data Traceability and Provenance}
\label{NFR3}

\subsubsection{Open-Source Software}
\label{NFR4}

\subsubsection{Reusability}
\label{NFR5}

\subsubsection{Usability for non-developers}
\label{NFR6}

\subsubsection{Configurability and Adaptability}
\label{NFR7}

\section{Conceptual Architecture}

\subsection{Pipeline Model}

\subsubsection{Pipeline}

\subsubsection{Step}

\subsubsection{Schedule}

\subsubsection{Data Set}

\subsubsection{Configurations}

\subsection{Domain Model}

\subsubsection{Project}

\subsubsection{Organization}

\subsubsection{Technology}

\subsubsection{Programme}

\subsubsection{Grant}

\section{Summary}
